\section{引言}

从视觉信息中进行3D目标检测（3D Object Detection）一直是低成本自动驾驶系统面临的长期挑战。虽然利用激光雷达（LiDAR）等传感器采集的点云进行目标检测得益于可见物体3D结构信息，但基于相机的设置更为不适定，因为必须仅从RGB图像中包含的2D信息生成3D边界框预测。

现有方法~\cite{zhou2019centernet,wang2021fcos3d}通常完全基于2D计算构建检测流程。它们使用为2D任务设计的检测流程（如CenterNet~\cite{zhou2019centernet}、FCOS~\cite{tian2019fcos}）来预测物体姿态和速度等3D信息，而未考虑3D场景结构或传感器配置。这些方法需要多个后处理步骤来融合不同相机的预测并去除冗余边界框，导致效率和效果之间存在严重权衡。

作为这些基于2D方法的替代方案，一些方法通过应用3D重建方法~\cite{monodepth2,bts,packnet}，从相机图像创建场景的伪激光雷达（Pseudo-LiDAR）或距离输入，从而将更多3D计算引入检测流程。然后，可对这些数据应用3D目标检测方法，就像直接来自3D传感器一样。然而，这种策略会受到累积误差的影响~\cite{demistifying}：估计不准确的深度值会对3D目标检测性能产生强烈的负面影响，而3D目标检测本身也可能产生误差。

本文提出了一种更优雅的2D观测与3D预测之间的转换方案，该方案不依赖密集深度预测模块。本文框架称为\name（多视角3D检测），以自顶向下的方式解决该问题。通过几何反向投影（Geometric Back-Projection）和相机变换矩阵，将2D特征提取与3D目标预测相连接。

本文方法从一组稀疏的目标先验开始，这些先验在数据集上共享并以端到端方式学习。为获取场景特定信息，将从这些目标先验解码得到的一组参考点反向投影到每个相机，并获取由ResNet骨干网络~\cite{He2015}提取的相应图像特征。参考点的图像特征随后通过多头自注意力层~\cite{Vaswani2017}进行交互。经过一系列自注意力层后，从每一层读取边界框参数，并使用受DETR~\cite{detr}启发的集合到集合损失来评估性能。

本文架构不进行点云重建或从图像显式预测深度，因此对深度估计误差具有鲁棒性。此外，本文方法不需要任何后处理（如非极大值抑制），提高了效率并减少了对人工设计的输出清洗方法的依赖。在nuScenes数据集上，本文方法（无NMS）与先前方法（有NMS）性能相当。在相机重叠区域，本文方法显著优于其他方法。

\textbf{贡献。} 本文的主要贡献总结如下：
\begin{itemize}
  \item 提出了一种基于RGB图像的简洁3D目标检测模型。与现有在最后阶段才融合不同相机视图目标预测的方法不同，本文方法在每一层计算中都融合了所有相机视图的信息。据本文所知，这是首次将多相机检测视为3D集合到集合预测的尝试。
  \item 引入了一个通过几何反向投影连接2D特征提取和3D边界框预测的模块。该模块不受辅助网络不准确深度预测的影响，并通过将3D信息反向投影到所有可用帧上无缝利用多相机信息。
  \item 与Object DGCNN~\cite{objdgcnn}类似，本文方法不需要逐图像或全局NMS等后处理，性能与现有基于NMS的方法相当。在相机重叠区域，本文方法以显著优势优于其他方法。
  \item 公开代码以促进可重复性和未来研究。\footnote{\url{https://github.com/WangYueFt/detr3d}}
\end{itemize}
